\documentclass[11pt]{article}
\usepackage[margin=1in]{geometry}
\usepackage{amsmath,amssymb,amsthm}
\usepackage{graphicx}
\usepackage[hidelinks]{hyperref}

\newtheorem{theorem}{Theorem}
\newtheorem{remark}{Remark}

\newcommand{\eps}{\varepsilon}
\newcommand{\ellone}{\ell_1}
\newcommand{\B}{\mathcal B}
\newcommand{\Sn}{S_n}
\newcommand{\dist}{d}

\title{A finite-cotype obstruction for the all-radius entropy version}
\author{Run fa\_banach\_001, lane 3}
\date{2026-06-15}

\begin{document}
\maketitle

\paragraph{Status.}
Candidate counterexample, likely valid after human review.  The result below
refutes the finite-cotype converse for the ``all $\eps>0$'' interpretation of
Pisier's problem on the entropy of $\Sn(X)$.  It does not solve the
``some $\eps>0$'' version.

\paragraph{Source problem.}
Pisier observes that if
\[
   \log\log N(\Sn(X),\dist,1+\eps)=o(n)
\]
for some $\eps>0$, then $X$ does not contain $\ell_\infty^n$ uniformly and
hence has finite cotype.  He then asks whether the converse holds and asks to
characterize the infinite-dimensional Banach spaces satisfying the displayed
condition ``for some (or all) $\eps>0$'' \cite[p. 9]{Pisier}.  The exact source
passage is also TeX lines 570--585 of the parsed source.

\begin{center}
\includegraphics[width=0.96\linewidth]{figures/open_problem_crop.png}
\end{center}

\paragraph{Source theorem used.}
We use Pisier's local entropy theorem \cite[Theorem 2.1]{Pisier}: for all
$1<r<R$ there is $c=c(r,R)>0$ such that, for every sufficiently large $n$ and
every $n$-dimensional Banach space $E$, there is a set
$T_n\subset \mathcal B_n$ with
\[
  |T_n|\ge \exp\exp(cn), \qquad
  \dist(t,u)>r \ (t\ne u), \qquad
  \dist(t,E)<R \ (t\in T_n).
\]
In the parsed source this is Theorem \texttt{local}, TeX lines 606--614.

\begin{theorem}
For every $s>1$ there is an infinite-dimensional Banach space $X_s$,
isomorphic to $\ell_1$, and a constant $a_s>0$ such that
\[
  \liminf_{n\to\infty} \frac{1}{n}
  \log\log N(\Sn(X_s),\dist,s) \ge a_s .
\]
Consequently, for every fixed $\eps_0>0$, with $s=1+\eps_0$, the space
$X_s$ has finite cotype but fails Pisier's entropy condition at the radius
$1+\eps_0$.
\end{theorem}

\paragraph{Proof intuition.}
The local theorem says that every finite-dimensional space, in particular
$\ell_1^n$, has a whole double-exponential cloud of $n$-dimensional spaces
nearby, separated at any prescribed smaller scale.  Put all those nearby
spaces directly into an $\ell_1$-sum.  Because every block is uniformly
isomorphic to $\ell_1^n$, the whole sum is still just an equivalent renorming
of $\ell_1$.  But at dimension $n$ the subspace set $\Sn(X_s)$ now contains
the entire separated cloud.

\begin{proof}
Fix $s>1$.  Choose $R>s^2$.  Apply the local theorem with $r=s^2$ and center
$E=\ell_1^n$.  There are constants $c>0$ and $n_0$ such that for every
$n\ge n_0$ one can choose a finite family
\[
   \mathcal F_n=\{E_{n,\alpha}:\alpha\in A_n\}\subset \mathcal B_n
\]
satisfying
\[
   |A_n|\ge \exp\exp(cn),\qquad
   \dist(E_{n,\alpha},E_{n,\beta})>s^2\quad(\alpha\ne\beta),
\]
and
\[
   \dist(E_{n,\alpha},\ell_1^n)<R\quad(\alpha\in A_n).
\]
After increasing $R$ slightly, choose for each $(n,\alpha)$ an isomorphism
$U_{n,\alpha}:E_{n,\alpha}\to \ell_1^n$ with
\[
   \|U_{n,\alpha}\|\,\|U_{n,\alpha}^{-1}\|\le R .
\]
Define
\[
   X_s=\left(\bigoplus_{n\ge n_0}\ \bigoplus_{\alpha\in A_n}
        E_{n,\alpha}\right)_{\ell_1}.
\]
The block diagonal map $\bigoplus_1 U_{n,\alpha}$ is an isomorphism from
$X_s$ onto
\[
   \left(\bigoplus_{n\ge n_0}\ \bigoplus_{\alpha\in A_n}
        \ell_1^n\right)_{\ell_1},
\]
with distortion at most $R$.  The latter space is isometric to $\ell_1$ on a
countable index set.  Therefore $X_s$ is isomorphic to $\ell_1$, and in
particular has cotype $2$.

For each $n\ge n_0$, every $E_{n,\alpha}$ is a coordinate subspace of $X_s$.
Thus $\mathcal F_n\subset \Sn(X_s)$.  If a $\dist$-ball of radius $s$ contained
two distinct members $E_{n,\alpha}$ and $E_{n,\beta}$, the multiplicative
triangle inequality for the Banach-Mazur distance would give
\[
   \dist(E_{n,\alpha},E_{n,\beta})\le s^2,
\]
contrary to the construction.  Hence any radius-$s$ cover of $\Sn(X_s)$ must
use at least one ball for each member of $\mathcal F_n$, and so
\[
   N(\Sn(X_s),\dist,s)\ge |\mathcal F_n|\ge \exp\exp(cn)
\]
for all $n\ge n_0$.  Therefore
\[
   \liminf_{n\to\infty} \frac{1}{n}
   \log\log N(\Sn(X_s),\dist,s) \ge c.
\]
This proves the theorem, with $a_s=c$.
\end{proof}

\begin{remark}[Why this is not the some-radius counterexample]
Because $X_s$ is isomorphic to $\ell_1$, the usual
Bourgain-Lindenstrauss-Milman embedding theorem cited by Pisier transfers
through the global isomorphism and gives the same $o(n)$ double-logarithmic
entropy estimate at sufficiently large covering radii.  Thus the example is
intentionally scoped: it disproves the all-$\eps$ finite-cotype converse and
shows that the all-radius property is not stable under equivalent renorming,
but it does not disprove the some-$\eps$ version of the characterization
problem.
\end{remark}

\paragraph{Verifier report.}
No computational verifier is relevant here.  The proof has three check points:
(1) Pisier's local theorem is applied with $r=s^2<R$ and center $\ell_1^n$;
(2) the $\ell_1$-sum of uniformly $R$-isomorphic finite-dimensional
$\ell_1$ blocks is $R$-isomorphic to $\ell_1$; and (3) separation $>s^2$
forces one radius-$s$ covering ball per selected block by the multiplicative
triangle inequality.  The scope caveat above also checks the apparent tension
with the BLM theorem for $\ell_1$.

\paragraph{Novelty check.}
The run's lightweight indexes were searched for \texttt{1306.5325},
\texttt{S\_n(X)}, \texttt{SQ\_n(X)}, \texttt{Banach-Mazur compactum},
\texttt{metric entropy}, and \texttt{finite cotype}.  No duplicate packet or
attempt was found.  External exact-phrase searches for the source problem did
not reveal a known answer.  Novelty confidence is moderate pending human
review.

\begin{thebibliography}{9}
\bibitem{Pisier}
Gilles Pisier,
\emph{On the metric entropy of the Banach-Mazur compactum},
arXiv:1306.5325.
\end{thebibliography}

\end{document}
